% Options for packages loaded elsewhere
\PassOptionsToPackage{unicode}{hyperref}
\PassOptionsToPackage{hyphens}{url}
\documentclass[
  ignorenonframetext,
  aspectratio=169,
]{beamer}
\newif\ifbibliography
\usepackage{pgfpages}
\setbeamertemplate{caption}[numbered]
\setbeamertemplate{caption label separator}{: }
\setbeamercolor{caption name}{fg=normal text.fg}
\beamertemplatenavigationsymbolsempty
% remove section numbering
\setbeamertemplate{part page}{
  \centering
  \begin{beamercolorbox}[sep=16pt,center]{part title}
    \usebeamerfont{part title}\insertpart\par
  \end{beamercolorbox}
}
\setbeamertemplate{section page}{
  \centering
  \begin{beamercolorbox}[sep=12pt,center]{section title}
    \usebeamerfont{section title}\insertsection\par
  \end{beamercolorbox}
}
\setbeamertemplate{subsection page}{
  \centering
  \begin{beamercolorbox}[sep=8pt,center]{subsection title}
    \usebeamerfont{subsection title}\insertsubsection\par
  \end{beamercolorbox}
}
% Prevent slide breaks in the middle of a paragraph
\widowpenalties 1 10000
\raggedbottom
\AtBeginPart{
  \frame{\partpage}
}
\AtBeginSection{
  \ifbibliography
  \else
    \frame{\sectionpage}
  \fi
}
\AtBeginSubsection{
  \frame{\subsectionpage}
}
\usepackage{iftex}
\ifPDFTeX
  \usepackage[T1]{fontenc}
  \usepackage[utf8]{inputenc}
  \usepackage{textcomp} % provide euro and other symbols
\else % if luatex or xetex
  \usepackage{unicode-math} % this also loads fontspec
  \defaultfontfeatures{Scale=MatchLowercase}
  \defaultfontfeatures[\rmfamily]{Ligatures=TeX,Scale=1}
\fi
\usepackage{lmodern}
\ifPDFTeX\else
  % xetex/luatex font selection
\fi
% Use upquote if available, for straight quotes in verbatim environments
\IfFileExists{upquote.sty}{\usepackage{upquote}}{}
\IfFileExists{microtype.sty}{% use microtype if available
  \usepackage[]{microtype}
  \UseMicrotypeSet[protrusion]{basicmath} % disable protrusion for tt fonts
}{}
\makeatletter
\@ifundefined{KOMAClassName}{% if non-KOMA class
  \IfFileExists{parskip.sty}{%
    \usepackage{parskip}
  }{% else
    \setlength{\parindent}{0pt}
    \setlength{\parskip}{6pt plus 2pt minus 1pt}}
}{% if KOMA class
  \KOMAoptions{parskip=half}}
\makeatother
\usepackage{longtable,booktabs,array}
\usepackage{caption}
\captionsetup[table]{skip=6pt}
\newcounter{none} % for unnumbered tables
\usepackage{calc} % for calculating minipage widths
\usepackage{caption}
% Make caption package work with longtable
\makeatletter
\def\fnum@table{\tablename~\thetable}
\makeatother
\setlength{\emergencystretch}{3em} % prevent overfull lines
\providecommand{\tightlist}{%
  \setlength{\itemsep}{0pt}\setlength{\parskip}{0pt}}
% Auto-generated by infrastructure/rendering/_pdf_latex_helpers.py::extract_math_font_preamble.
% Minimal header passed to Pandoc via -H so Beamer slides pick up the same math font
% as the combined-PDF path. Adjust the manuscript's preamble.md to change the font globally.
\usepackage{unicode-math}
\setmathfont{latinmodern-math.otf}

% Manuscript macro/package fallbacks (newcommand -> providecommand).
% Core mathematics
\usepackage{amsmath}
\usepackage{amssymb}
% Document layout
% Tables
\usepackage{booktabs}
% Caption styling (booktabs already loaded above). Small caption font with a
% bold label ("Figure 1", "Table 2") gives the math/figure-heavy layout a
% consistent, journal-like caption rhythm without fighting pandoc-crossref —
% pandoc emits standard \caption{}, and caption/captionsetup only restyle it.
% ── Theorem environments ─────────────────────────────────────────────────
% amsthm is loaded above. The FedGVI<->Friston bridge is stated as a small
% number of theorems/definitions/lemmas/propositions/corollaries; sharing one
% counter (definition/lemma/proposition/corollary all step the `theorem`
% counter) gives a single monotone numbering 1, 2, 3, ... across all kinds, so
% authors can reference them as Theorem 1, Definition 2, Corollary 3 without a
% per-kind counter drifting out of order. Plain style for theorem-like results,
% definition style (upright body) for definitions.
% (algorithm2e intentionally not loaded: this manuscript states results as
% numbered amsthm Theorems/Definitions, not algorithm2e pseudocode.)
% Code listings
\usepackage{listings}
% Typography and formatting
% (siunitx intentionally not loaded: no \SI/\num units appear in this manuscript.)
% Cross-references and citations
% (cleveref and titlesec intentionally not loaded: unavailable in this TeX tree
% and not required. Cross-references resolve through pandoc-crossref's own
% \ref-style names — no \cref appears — and section heading styling falls back
% to the pandoc/LaTeX defaults.)
% ── TeX Live Basic-compatible mono font for code listings ────────────
% Latin Modern Mono is bundled with TeX Live Basic and keeps the manuscript
% renderable on minimal TeX installations. Mathematical Unicode coverage is
% provided separately by Latin Modern Math below, so the code font does not
% need a private JuliaMono installation.
%
% Use the TeX font filename rather than the system-family name: the minimal
% TeX Live fontconfig database may not register the OpenType family even though
% kpsewhich can resolve the bundled files.
% Math font for unicode-math: Latin Modern Math (TeX Live) has full BMP
% coverage including U+2223 (\mid), U+226A/226B (\ll/\gg), and the Greek/
% blackboard letters used in equations. Without an explicit \setmathfont,
% unicode-math falls back to lmroman text font which lacks several glyphs
% and emits "Missing character" warnings on every \mid in math mode.

% Slide layout defaults for warning-clean scientific decks.
\usepackage{etoolbox}
\IfFileExists{xurl.sty}{\usepackage{xurl}}{}
\IfFileExists{seqsplit.sty}{\usepackage{seqsplit}}{\newcommand{\seqsplit}[1]{#1}}
\protected\def\breakseq#1{\seqsplit{#1}}
\protected\def\breaktt#1{\begingroup\ttfamily\seqsplit{#1}\endgroup}
\setlength{\emergencystretch}{6em}
\tolerance=5000
\hbadness=10000
\hfuzz=1pt
\setlength{\tabcolsep}{2pt}
\AtBeginEnvironment{longtable}{\tiny\renewcommand{\arraystretch}{0.86}\setlength{\tabcolsep}{1pt}}
\AtBeginEnvironment{tabular}{\tiny\renewcommand{\arraystretch}{0.86}\setlength{\tabcolsep}{1pt}}
\AtBeginEnvironment{equation}{\tiny}
\AtBeginEnvironment{equation*}{\tiny}
\AtBeginEnvironment{align}{\tiny}
\AtBeginEnvironment{align*}{\tiny}
\AtBeginEnvironment{itemize}{\footnotesize}
\AtBeginEnvironment{enumerate}{\footnotesize}
\AtBeginEnvironment{description}{\footnotesize}
\setbeamerfont{caption}{size=\tiny}
\setbeamerfont{caption name}{size=\tiny}
\setbeamerfont{normal text}{size=\small}
\setbeamerfont{frametitle}{size=\small}
\setbeamerfont{section title}{size=\footnotesize}
\setbeamerfont{subsection title}{size=\footnotesize}
\setbeamertemplate{section page}{%
  \centering
  \begin{beamercolorbox}[sep=12pt,center,wd=\paperwidth]{section title}
    \parbox{0.86\paperwidth}{\centering\usebeamerfont{section title}\insertsection\par}
  \end{beamercolorbox}
}
\setbeamertemplate{subsection page}{%
  \centering
  \begin{beamercolorbox}[sep=8pt,center,wd=\paperwidth]{subsection title}
    \parbox{0.86\paperwidth}{\centering\usebeamerfont{subsection title}\insertsubsection\par}
  \end{beamercolorbox}
}
\setlength{\abovecaptionskip}{2pt}
\setlength{\belowcaptionskip}{0pt}

% Natbib and cross-reference fallbacks — slides don't load natbib
% or cleveref, but combined-PDF manuscript prose may emit these
% commands. The fallback renders citations as a bracketed key list
% and cross-references as detokenized labels so slides don't fail on
% undefined control sequences or raw underscores. \providecommand is
% a no-op if packages load later.
\providecommand{\citep}[1]{[#1]}
\providecommand{\citet}[1]{#1}
\providecommand{\citealp}[1]{#1}
\providecommand{\citeauthor}[1]{#1}
\providecommand{\citeyear}[1]{#1}
\providecommand{\cref}[1]{\texttt{\detokenize{#1}}}
\providecommand{\Cref}[1]{\texttt{\detokenize{#1}}}
\providecommand{\crefrange}[2]{\texttt{\detokenize{#1}}--\texttt{\detokenize{#2}}}
\providecommand{\Crefrange}[2]{\texttt{\detokenize{#1}}--\texttt{\detokenize{#2}}}
\providecommand{\paragraph}[1]{\textbf{#1}\ }

% Opt-in accessible presentation profile.
\setbeamerfont{normal text}{size*={20pt}{24pt}}
\setbeamerfont{frametitle}{size*={28pt}{32pt}}
\setbeamerfont{section title}{size*={28pt}{32pt}}
\setbeamerfont{subsection title}{size*={28pt}{32pt}}
\setbeamerfont{caption}{size*={16pt}{19pt}}
\setbeamerfont{caption name}{size*={16pt}{19pt}}
\setbeamerfont{itemize/enumerate body}{size*={20pt}{24pt}}
\setbeamerfont{itemize/enumerate subbody}{size*={20pt}{24pt}}
\setbeamerfont{itemize/enumerate subsubbody}{size*={20pt}{24pt}}
\setbeamerfont{itemize item}{size*={20pt}{24pt}}
\setbeamerfont{itemize subitem}{size*={20pt}{24pt}}
\setbeamerfont{itemize subsubitem}{size*={20pt}{24pt}}
\setbeamerfont{description body}{size*={20pt}{24pt}}
\setbeamerfont{description item}{size*={20pt}{24pt}}
\setbeamerfont{quote}{size*={20pt}{24pt}}
\setbeamerfont{quotation}{size*={20pt}{24pt}}
\AtBeginEnvironment{longtable}{\fontsize{20pt}{24pt}\selectfont}
\AtBeginEnvironment{tabular}{\fontsize{20pt}{24pt}\selectfont}
\AtBeginEnvironment{equation}{\fontsize{20pt}{24pt}\selectfont}
\AtBeginEnvironment{equation*}{\fontsize{20pt}{24pt}\selectfont}
\AtBeginEnvironment{align}{\fontsize{20pt}{24pt}\selectfont}
\AtBeginEnvironment{align*}{\fontsize{20pt}{24pt}\selectfont}
\AtBeginEnvironment{itemize}{\fontsize{20pt}{24pt}\selectfont}
\AtBeginEnvironment{enumerate}{\fontsize{20pt}{24pt}\selectfont}
\AtBeginEnvironment{description}{\fontsize{20pt}{24pt}\selectfont}
\AtBeginDocument{\usebeamerfont{normal text}}
\newlength{\AFSlideTableWidth}
% The fixed footer reservation is calibrated with the semantic seven-line regular-body budget.
\setbeamertemplate{footline}{%
  \leavevmode\vbox{%
    \hbox{%
      \begin{beamercolorbox}[wd=\paperwidth,ht=16pt,dp=3pt,center]{author in head/foot}%
        \fontsize{16pt}{19pt}\selectfont Untagged PDF derivative \textbar\ \href{../web/index.html}{HTML reader}%
      \end{beamercolorbox}%
    }%
    \vskip6pt%
  }%
}
% Accessible footer reserved height: 25pt.

\newtheorem{proposition}{Proposition}
\newtheorem{hypothesis}{Hypothesis}
\newtheorem{remark}[theorem]{Remark}
\newtheorem{axiom}{Axiom}
\newtheorem{property}{Property}
\usepackage{bookmark}
\IfFileExists{xurl.sty}{\usepackage{xurl}}{} % add URL line breaks if available
\urlstyle{same}
% fallback for those not using the hyperref driver hyperxmp:
\makeatletter
\@ifundefined{xmpquote}{\newcommand{\xmpquote}[1]{#1}}{}
\makeatother
\hypersetup{
  hidelinks,
  pdfcreator={LaTeX via pandoc}}

\author{\texorpdfstring{}{}}
\date{}

\begin{document}

\begin{frame}{Learning stack: EFE, Dirichlet updates, and BMR}
\protect\phantomsection\label{sec:methods-learning}
Active-inference agents do more than infer states under a fixed model:
they learn the parameters of the model, score policies by expected free
energy, and revise model structure.
\end{frame}

\begin{frame}{Learning stack (part 2)}
The active-inference community has standardized all three operations (Da
Costa et al. 2020; Smith et al. 2022), and Friston et al. (2024) place
them at the heart of the federated belief-sharing scenario.
\end{frame}

\begin{frame}{Learning stack (part 3)}
We reimplement each in closed form so the language-acquisition,
expected-free-energy, and emergence studies of
sec.~5.22 rest on machine-checkable
quantities rather than fitted curves.
\end{frame}

\begin{frame}{Conjugate Dirichlet learning from co-occurrence counts}
\protect\phantomsection\label{sec:methods-dirichlet}
A sentinel learns its observation model \(A\) by placing a Dirichlet
prior with concentration \(a\) on each column and updating it
conjugately from observation-state co-occurrence counts (Friston et al.
(2024), their equations 9--12).
\end{frame}

\begin{frame}{Conjugate Dirichlet learning\ldots{} (part 2)}
One learning step adds the expected sufficient statistics for that step
and reads off the column-normalized posterior mean,
\end{frame}

\begin{frame}{Conjugate Dirichlet learning\ldots{} (part 3)}
\begin{equation}\tag{18}\protect\phantomsection\label{eq:dirichlet-update}{
a \;\leftarrow\; a + \text{counts},
\qquad
\mathbb{E}[A]_{o s} \;=\; \frac{a_{o s}}{\sum_{o'} a_{o' s}},
}\end{equation}
\end{frame}

\begin{frame}[fragile]{Conjugate Dirichlet learning\ldots{} (part 4)}
implemented by \texttt{learn\_likelihood} in the
\texttt{dirichlet\_learning} module. Intuitively, each concentration
vector is a running tally of how often each outcome was seen while the
creature occupied a given state: the prior seeds that tally with
pseudo-counts, every step adds the co-occurrences it witnessed,
\end{frame}

\begin{frame}{Conjugate Dirichlet learning\ldots{} (part 5)}
and the posterior mean is simply the tally renormalized into a
categorical.

Likelihood learning is therefore bookkeeping --- accumulate counts, then
normalize --- with no iterative optimization loop.
\end{frame}

\begin{frame}{Conjugate Dirichlet learning\ldots{} (part 6)}
We drive eq.~18 with the expected sufficient
statistics under the true model --- a fixed count batch
\(\text{count\_scale}\cdot A^{\star}\) per step, optionally jittered by
a seeded generator. As the concentrations accumulate,
\end{frame}

\begin{frame}{Conjugate Dirichlet learning\ldots{} (part 7)}
the expected likelihood \(\mathbb{E}[A]\) converges to the
data-generating \(A^{\star}\); convergence is measured by the per-column
KL divergence summed over hidden states,
\end{frame}

\begin{frame}{Conjugate Dirichlet learning\ldots{} (part 8)}
\begin{equation}\tag{19}\protect\phantomsection\label{eq:dirichlet-kl}{
\begin{aligned}
\mathrm{KL}\big(A^{\star}\|\mathbb{E}[A]\big)
&= \sum_s\sum_o A^{\star}_{os}\\
&\quad\cdot\ln\frac{A^{\star}_{os}}{\mathbb{E}[A]_{os}}.
\end{aligned}
}\end{equation}
\end{frame}

\begin{frame}{Conjugate Dirichlet learning\ldots{} (part 9)}
which decreases monotonically toward zero --- the standard-Bayes / KL
fixed point. The learned likelihood always has full support (the
Dirichlet prior is strictly positive), so eq.~19 is
finite throughout. ISC-17 pins the monotone-decreasing KL trajectory,
\end{frame}

\begin{frame}{Conjugate Dirichlet learning\ldots{} (part 10)}
and the language-acquisition study of
sec.~5.22 reports the descent of
eq.~19 across 24 steps.
\end{frame}

\begin{frame}{Conjugate Dirichlet learning\ldots{} (part 11)}
The implementation also carries the \(\eta\) forgetting hyperprior of
Friston et al. (2024), their equation 12: before each conjugate addition
the running concentration is decayed so the \emph{total} concentration
mass saturates at \(\eta\) rather than growing without bound,
\end{frame}

\begin{frame}{Conjugate Dirichlet learning\ldots{} (part 12)}
modeling an agent that stays adaptable instead of becoming infinitely
confident. With \(\eta\) unset the classical unbounded accumulation of
eq.~18 is recovered.
\end{frame}

\begin{frame}{Expected free energy as the action-selection objective}
\protect\phantomsection\label{sec:methods-efe}
A sentinel scores a candidate policy \(\boldsymbol{\pi}\) by its
expected free energy \(G(\boldsymbol{\pi})\), which the active-inference
formulation decomposes two equivalent ways (Friston et al. (2024), their
equation 2):
\end{frame}

\begin{frame}{Expected free energy as\ldots{} (part 2)}
a cost view of risk plus ambiguity, and a value view of pragmatic plus
epistemic value.
\end{frame}

\begin{frame}{Expected free energy as\ldots{} (part 3)}
The two views are the same scalar rearranged, stated as
eq.~22 and pinned to a zero residual by the
algebraic identity eq.~23 in
sec.~6.
\end{frame}

\begin{frame}[fragile]{Expected free energy as\ldots{} (part 4)}
We compute every term in closed form from the categorical model
\((A, B, C, D_0)\) in \texttt{expected\_free\_energy.decompose}:
\end{frame}

\begin{frame}{Risk}
\protect\phantomsection\label{risk}
\emph{Risk} is
\(\mathrm{KL}\big(q(o\mid\boldsymbol{\pi})\,\|\,p_C(o)\big)\), the
deviation of the policy-predicted outcomes from the preferred-outcome
pmf \(p_C(o)=\mathrm{softmax}(C)[o]\). Write
\(q_{\boldsymbol{\pi}}(o):=q(o\mid\boldsymbol{\pi})\) for this
scored-policy outcome predictive,
\end{frame}

\begin{frame}{Risk (part 2)}
used in the remaining terms.
\end{frame}

\begin{frame}{Ambiguity}
\protect\phantomsection\label{ambiguity}
\emph{Ambiguity} is the expected likelihood entropy
\(\mathbb{E}_{q(s)}\!\big[H[p(o\mid s)]\big]\), the outcome uncertainty
given the state.
\end{frame}

\begin{frame}{Pragmatic value}
\protect\phantomsection\label{pragmatic-value}
\emph{Pragmatic value} is the expected log-preference
\(\mathbb{E}_{q_{\boldsymbol{\pi}}(o)}[\ln p_C(o)]\) --- the utility,
exploitation term.
\end{frame}

\begin{frame}{Epistemic value}
\protect\phantomsection\label{epistemic-value}
\emph{Epistemic value} is the state-outcome mutual information
\(H[q_{\boldsymbol{\pi}}(o)] - \mathbb{E}_{q(s)}[H[p(o\mid s)]]\) ---
the expected information gain that drives exploration.
\end{frame}

\begin{frame}[fragile]{Epistemic value (part 2)}
Because there is no sampling, the identity of eq.~23
holds to floating-point tolerance; ISC-19
(\texttt{expected\_free\_energy}) pins the residual of the decomposition
to zero and pins each term's semantics independently (deterministic
likelihoods give zero ambiguity;
\end{frame}

\begin{frame}{Epistemic value (part 3)}
uninformative likelihoods give zero epistemic value; preference-matched
predictions lower risk).
\end{frame}

\begin{frame}{Epistemic value (part 4)}
fig.~8 visualizes the additive cost view and the
signed pragmatic/epistemic waterfall terminating at
\(G(\boldsymbol{\pi})\) --- a deterministic identity (Proposition
7), not a fitted result.
\end{frame}

\begin{frame}{Bayesian model reduction for the configured pruning sign
control}
\protect\phantomsection\label{sec:methods-bmr}
Sentinels also revise model \emph{structure}. Bayesian model reduction
(BMR) scores whether a reduced model --- for example one that prunes a
redundant location column by shrinking its concentration toward zero ---
\end{frame}

\begin{frame}{Bayesian model reduction\ldots{} (part 2)}
has more evidence than the full model, \emph{without re-running
inference} (Friston \& Penny via the post-hoc model optimization lineage
(Friston and Penny 2011); the same Beta-function identity is their
equation 13 (Friston et al. (2024))).
\end{frame}

\begin{frame}{Bayesian model reduction\ldots{} (part 3)}
Because the likelihood is shared, the reduced posterior is available in
closed form,
\(\text{reduced\_post} = \text{post} + \text{reduced\_prior} - \text{prior}\),
and the change in (negative) variational free energy is a difference of
log multivariate Beta functions,
\end{frame}

\begin{frame}{Bayesian model reduction\ldots{} (part 4)}
\begin{equation}\tag{20}\protect\phantomsection\label{eq:bmr-deltaf}{
\begin{aligned}
\Delta F
&\;=\; \ln B(\text{prior}) - \ln B(\text{post})\\
&\qquad + \ln B(\text{reduced\_post})\\
&\qquad - \ln B(\text{reduced\_prior}),\\
\ln B(a)
&\;=\; \textstyle\sum_k \ln\Gamma(a_k)
 - \ln\Gamma\!\big(\sum_k a_k\big).
\end{aligned}
}\end{equation}
\end{frame}

\begin{frame}[fragile]{Bayesian model reduction\ldots{} (part 5)}
computed in \texttt{bayesian\_model\_reduction.reduce}. A positive
\(\Delta F\) in eq.~20 means the reduced model carries
more evidence --- the pruned structure was redundant and should be
adopted; a negative \(\Delta F\) means the reduction destroyed support
the data require.
\end{frame}

\begin{frame}{Bayesian model reduction\ldots{} (part 6)}
When the reduced prior equals the prior the score is identically zero in
exact algebra, a zero point the suite pins to machine precision
(ISC-20).
\end{frame}

\begin{frame}{Bayesian model reduction\ldots{} (part 7)}
The configured BMR sign-control study of
sec.~5.22 uses this operation over
\(n = 4\) candidate states. It contrasts a redundant reduction
(\(\Delta F = 3.68\) nats; adopt) with a supported one
(\(\Delta F = -27.67\) nats; reject) in fig.~12.
\end{frame}

\begin{frame}{Bayesian model reduction\ldots{} (part 8)}
This fixed-candidate algebraic comparison is deterministic, so it has no
resampled sample or bootstrap interval and does not establish structure
discovery or universal emergence.
\end{frame}

\end{document}
